\chapter{Phonemic inventory}

\section{Vowels}

\subsection{Monophthongs}

Table~\vref{tab:vowel-phonemes} shows Interturkic's eight vowel phonemes,
arranged according to their height, backness, and rounding.
In native vocabulary, low rounded vowels may occur only in the first syllables
of roots.
All other vowels are unrestriced in their distribution, and may occur in any
syllable of any word.

\begin{table}
  \centering
  \caption{Vowel phonemes}
  \label{tab:vowel-phonemes}
  \begin{tabular}{lllll}
    \toprule
    \multirow{2}{*}{} & \multicolumn{2}{l}{Front} & \multicolumn{2}{l}{Back} \\
    \cmidrule(lr){2-3} \cmidrule(lr){4-5}
    & unrounded & rounded & unrounded & rounded \\
    \midrule
    High & i & y & ɯ & u \\
    Low & ɛ & œ & ɑ & ɔ \\
    \bottomrule
  \end{tabular}
\end{table}

\subsection{Diphthongs}

Interturkic has a small set of diphthongs, formed from monophthongs followed by
glides, as in (\vref{ex:diphthong-yes}).
Note that such sequences are \emph{not} considered diphthongs if split across a
syllable boundary, as in (\vref{ex:diphthong-no}).
In any case, diphthongs are still analysed as \emph{pairs} of phonemes, rather
than as phonemes in and of themselves.

\pex
  \a\label{ex:diphthong-yes}%
    \vtop{\halign{%
      \phonem{#}\hfil\enspace&`#'\cr
      ɑj&moon, month\cr
      tɑw&mountain\cr}}
  \a\label{ex:diphthong-no}%
    \vtop{\halign{%
      \phonem{#}\hfil\enspace&`#'\cr
      bɔ.jun&neck\cr
      ɑ.w&mountain\cr}}
\xe

Diphthongs formed with \phonem{w} are treated as rounded for the purposes of
vowel harmony.

\section{Consonants}